\chapter{The Transverse-Field XY Chain as a Mother Model}
\label{chap:xy-chain-mother-model}

The transverse-field XY chain is the first genuine mother model in these notes.
It is not merely one more spin chain in a list of examples.  It contains, as
parameter limits or equivalent fermionic descriptions, the transverse-field
Ising model, the XX chain, and the spinless Kitaev chain.  It is also the
simplest setting in which all technical tools developed in Part~I meet in one
calculation: Jordan--Wigner fermionization, parity-dependent boundary sectors,
Fourier transformation, BdG blocks, Bogoliubov quasiparticles, Gaussian
correlators, duality, and two-band topology.

The purpose of this chapter is therefore twofold.  First, we solve the XY chain
completely as a quadratic fermion problem.  Second, we organize its special
limits so that the later chapters can be read as controlled reductions of one
solved mother model.

\section{Hamiltonian and conventions}
\label{sec:xy-hamiltonian-conventions}

The transverse-field XY chain is defined by
\begin{equation}
\label{eq:xy-hamiltonian}
    H_{\rm XY}
    =-
    \sum_{j=1}^{L}
    \left[
        \frac{J(1+\gamma)}{2}\sigma_j^x\sigma_{j+1}^x
        +
        \frac{J(1-\gamma)}{2}\sigma_j^y\sigma_{j+1}^y
    \right]
    -h\sum_{j=1}^{L}\sigma_j^z .
\end{equation}
Here $J$ sets the exchange scale, $h$ is the transverse field, and $\gamma$ is the
XY anisotropy.  The isotropic XX chain is obtained at $\gamma=0$, while the
transverse-field Ising chain is obtained at $\gamma=1$.

Unless otherwise stated, the bulk formulas below are independent of whether one
starts from open or periodic spin boundary conditions.  Boundary conditions only
enter when the Jordan--Wigner string wraps around the chain.  This distinction is
central enough that we keep it explicit throughout the derivation.

The spin and fermion conventions are those fixed in Part~I:
\begin{equation}
\label{eq:xy-mother-jw-convention-reminder}
    c_j
    =
    \left(\prod_{\ell=1}^{j-1}\sigma_\ell^z\right)\sigma_j^+,
    \qquad
    c_j^\dagger
    =
    \left(\prod_{\ell=1}^{j-1}\sigma_\ell^z\right)\sigma_j^- ,
\end{equation}
with
\begin{equation}
    n_j=c_j^\dagger c_j=\frac{1-\sigma_j^z}{2},
    \qquad
    \sigma_j^z=1-2n_j.
\end{equation}
Thus a spin-down state is interpreted as an occupied fermion state.  This is the
convention responsible for the sign of the chemical potential in the Kitaev-chain
dictionary below.

It is useful to rewrite the exchange part of \cref{eq:xy-hamiltonian}
as
\begin{equation}
\label{eq:xy-mother-isotropic-anisotropic-split}
    \frac{J(1+\gamma)}{2}\sigma_j^x\sigma_{j+1}^x
    +
    \frac{J(1-\gamma)}{2}\sigma_j^y\sigma_{j+1}^y
    =
    \frac{J}{2}
    \left(
        \sigma_j^x\sigma_{j+1}^x
        +
        \sigma_j^y\sigma_{j+1}^y
    \right)
    +
    \frac{J\gamma}{2}
    \left(
        \sigma_j^x\sigma_{j+1}^x
        -
        \sigma_j^y\sigma_{j+1}^y
    \right).
\end{equation}
The first bracket becomes fermion hopping; the second becomes $p$-wave pairing.
This is the algebraic reason why the XY chain is a BdG mother model.

\section{Special limits and the mother-model viewpoint}
\label{sec:xy-special-limits-mother-model}

The most important reductions of \cref{eq:xy-hamiltonian} are
summarized in \cref{tab:xy-mother-special-limits}.

\begin{table}[h]
\centering
\small
\setlength{\tabcolsep}{3pt}
\caption{Special limits and equivalent descriptions of the transverse-field XY chain.}
\label{tab:xy-mother-special-limits}
\begin{tabularx}{\textwidth}{L{0.23\textwidth}L{0.29\textwidth}Y}
\toprule
Parameter choice or map & Model obtained & Interpretation \\
\midrule
$\gamma=1$ & Transverse-field Ising model & Maximal anisotropy; Ising criticality at $h=\pm J$ in the convention of this chapter \\
$\gamma=0$ & XX chain in a transverse field & Number-conserving free fermions; gapless for $\abs{h}<\abs{J}$ \\
$h=0$ & Zero-field anisotropic XY chain & Pure competition between hopping and pairing after Jordan--Wigner transformation \\
General $h,\gamma$ & Spinless BdG chain & Independent $(k,-k)$ two-level systems after Fourier transformation \\
Jordan--Wigner image & Kitaev $p$-wave chain & $t=J$, $\mu=-2h$, and pairing amplitude fixed by $J\gamma$ up to pairing convention \\
Near $h=\pm J$, $\gamma\neq0$ & Massive Majorana theory & Ising universality class with a sign-changing Majorana mass \\
Near $\gamma=0$, $\abs{h}<\abs{J}$ & Two Fermi points & XX/Luttinger-liquid universality class rather than Ising criticality \\
\bottomrule
\end{tabularx}
\end{table}

The role of the XY chain as a mother model can be summarized as
\begin{equation}
\label{eq:xy-mother-model-map}
\begin{aligned}
    \text{XY spin chain}
    &\xrightarrow{\rm JW}
    \text{spinless BdG chain}\\
    &\xrightarrow{\rm Fourier}
    \prod_{k>0}\text{two-level BdG blocks}
    \xrightarrow{\rm Bogoliubov}
    \text{free quasiparticles}.
\end{aligned}
\end{equation}
All exact results in the later TFIM, XX, and Kitaev-chain chapters are obtained
by specializing this chain of maps.

\section{Jordan--Wigner image}
\label{sec:xy-jordan-wigner-image}

We now derive the fermionic Hamiltonian.  The nearest-neighbor Jordan--Wigner
identities established earlier give
\begin{align}
\label{eq:xy-mother-jw-xx-plus-yy}
    \sigma_j^x\sigma_{j+1}^x+
    \sigma_j^y\sigma_{j+1}^y
    &=
    2\left(
        c_j^\dagger c_{j+1}+c_{j+1}^\dagger c_j
    \right),
    \\
\label{eq:xy-mother-jw-xx-minus-yy}
    \sigma_j^x\sigma_{j+1}^x-
    \sigma_j^y\sigma_{j+1}^y
    &=
    2\left(
        c_j^\dagger c_{j+1}^\dagger+c_{j+1}c_j
    \right),
    \\
\label{eq:xy-mother-jw-z}
    \sigma_j^z&=1-2c_j^\dagger c_j.
\end{align}
Substituting these identities into \cref{eq:xy-hamiltonian} gives the
bulk Jordan--Wigner image
\begin{equation}
\label{eq:xy-mother-bulk-fermion-hamiltonian}
    H_{\rm XY}^{\rm bulk}
    =
    -J\sum_j
    \left(
        c_j^\dagger c_{j+1}+c_{j+1}^\dagger c_j
    \right)
    -J\gamma\sum_j
    \left(
        c_j^\dagger c_{j+1}^\dagger+c_{j+1}c_j
    \right)
    -h\sum_j(1-2n_j).
\end{equation}
Equivalently,
\begin{equation}
\label{eq:xy-mother-bulk-fermion-expanded}
    H_{\rm XY}^{\rm bulk}
    =
    -hL
    +2h\sum_j n_j
    -J\sum_j
    \left(
        c_j^\dagger c_{j+1}+c_{j+1}^\dagger c_j
    \right)
    -J\gamma\sum_j
    \left(
        c_j^\dagger c_{j+1}^\dagger+c_{j+1}c_j
    \right).
\end{equation}
The field term therefore plays the role of a chemical-potential term for the
Jordan--Wigner fermions.

\begin{remark}[Why the anisotropy becomes pairing]
The isotropic part
$\sigma^x\sigma^x+\sigma^y\sigma^y$ conserves the number of spin flips and becomes
fermion hopping.  The anisotropic part
$\sigma^x\sigma^x-\sigma^y\sigma^y$ creates or annihilates spin flips in pairs and
therefore becomes a superconducting pairing term.  Fermion number is not
conserved for $\gamma\neq0$, but fermion parity is conserved.
\end{remark}

For an open chain, \cref{eq:xy-mother-bulk-fermion-hamiltonian} is the full
fermionic Hamiltonian, with the sums running over $j=1,\ldots,L-1$.

For a periodic spin chain, the boundary term carries the Jordan--Wigner parity
sign.  Let
\begin{equation}
\label{eq:xy-mother-fermion-parity}
    P_F=(-1)^N,
    \qquad
    N=\sum_{j=1}^{L}n_j,
\end{equation}
and fix a parity sector
\begin{equation}
    P_F=p,
    \qquad
    p=\pm1.
\end{equation}
As shown in the Jordan--Wigner chapter, the periodic spin chain is represented in
that sector by the effective fermionic boundary condition
\begin{equation}
\label{eq:xy-mother-effective-fermion-boundary}
    c_{L+1}=-p\,c_1.
\end{equation}
Thus
\begin{equation}
\label{eq:xy-mother-pbc-apbc-sector}
    p=+1
    \quad\Longleftrightarrow\quad
    c_{L+1}=-c_1
    \quad\text{(anti-periodic fermions)},
\end{equation}
whereas
\begin{equation}
\label{eq:xy-mother-pbc-pbc-sector}
    p=-1
    \quad\Longleftrightarrow\quad
    c_{L+1}=c_1
    \quad\text{(periodic fermions)}.
\end{equation}
In a fixed sector, one may use
\cref{eq:xy-mother-bulk-fermion-hamiltonian} with $j=1,\ldots,L$ provided the
boundary condition \cref{eq:xy-mother-effective-fermion-boundary} is imposed.  The
exact finite-size spin partition function or spectrum is obtained only after
combining the two parity sectors with the correct parity projection.  In the
thermodynamic bulk spectrum, this distinction affects only finite-size momentum
quantization.

\section{Momentum quantization and Fourier transform}
\label{sec:xy-momentum-quantization-fourier}

In a fixed parity sector $p$, the allowed momenta are
\begin{equation}
\label{eq:xy-mother-allowed-momenta}
    k_m=\frac{2\pi m+\vartheta_p}{L},
    \qquad
    m=0,1,\ldots,L-1,
\end{equation}
where
\begin{equation}
\label{eq:xy-mother-boundary-twist-sector}
    \vartheta_p=
    \begin{cases}
        \pi, & p=+1 \quad \text{(even sector, APBC)},\\
        0, & p=-1 \quad \text{(odd sector, PBC)}.
    \end{cases}
\end{equation}
We write the corresponding momentum set as $\mathcal{K}_{\vartheta_p}$.

We use the Fourier convention developed earlier,
\begin{equation}
\label{eq:xy-mother-fourier-convention}
    c_k
    =
    \frac{\ee^{-\ii\phi}}{\sqrt{L}}
    \sum_{j=1}^{L}\ee^{-\ii kj}c_j,
    \qquad
    c_j
    =
    \frac{\ee^{\ii\phi}}{\sqrt{L}}
    \sum_{k\in\mathcal{K}_{\vartheta_p}}\ee^{\ii kj}c_k .
\end{equation}
The twist $\vartheta_p$ is physical boundary data.  The constant phase $\phi$ is a
Fourier-gauge convention; it rotates the phase of the pairing gap but not the
spectrum.

The hopping and field terms give
\begin{equation}
\label{eq:xy-mother-normal-dispersion}
    -J\sum_j
    \left(
        c_j^\dagger c_{j+1}+c_{j+1}^\dagger c_j
    \right)
    +2h\sum_j n_j
    =
    \sum_{k\in\mathcal{K}_{\vartheta_p}}
    \xi(k)c_k^\dagger c_k,
\end{equation}
with
\begin{equation}
\label{eq:xy-mother-xi-k}
    \xi(k)=2(h-J\cos k).
\end{equation}
The pairing term gives an odd gap function.  With the real-space sign in
\cref{eq:xy-mother-bulk-fermion-hamiltonian}, the direct Fourier transform gives
\begin{equation}
\label{eq:xy-mother-raw-gap}
    \Delta_\phi^{\rm raw}(k)
    =
    -2\ii J\gamma\,\ee^{-2\ii\phi}\sin k.
\end{equation}
This sign is a convention for the orientation of the Nambu pairing plane.  Many
BdG summaries, including some compact formulas in these notes, use the equivalent
orientation
\begin{equation}
\label{eq:xy-mother-standard-gap}
    \Delta_\phi(k)
    =
    +2\ii J\gamma\,\ee^{-2\ii\phi}\sin k.
\end{equation}
The two choices are related by a harmless reversal of the pairing-plane
orientation, for example by $\gamma\mapsto-\gamma$ or by an equivalent Nambu-gauge
relabelling.  They have the same $\abs{\Delta(k)}$, the same quasiparticle
spectrum, the same critical lines, and the same absolute winding number.  Only
orientation-sensitive signs of winding numbers change.  In the rest of this
chapter we display formulas with the convention
\cref{eq:xy-mother-standard-gap}, because it matches the BdG-vector convention
used in the preceding chapters.

With this convention, the momentum-space Hamiltonian is
\begin{equation}
\label{eq:xy-mother-momentum-hamiltonian}
    H_{\rm XY}^{(p)}
    =
    -hL
    +
    \sum_{k\in\mathcal{K}_{\vartheta_p}}
    \xi(k)c_k^\dagger c_k
    +
    \frac12
    \sum_{k\in\mathcal{K}_{\vartheta_p}}
    \left[
        \Delta_\phi(k)c_k^\dagger c_{-k}^\dagger
        +
        \Delta_\phi(k)^*c_{-k}c_k
    \right],
\end{equation}
up to the finite-size parity projection already discussed.

\section{BdG block and quasiparticle spectrum}
\label{sec:xy-bdg-block-spectrum}

Introduce the Nambu spinor
\begin{equation}
\label{eq:xy-mother-nambu-spinor}
    \Psi_k=
    \begin{pmatrix}
        c_k\\ c_{-k}^\dagger
    \end{pmatrix}.
\end{equation}
For $k\not\equiv -k$ modulo $2\pi$, the pair $(k,-k)$ is governed by the BdG
block
\begin{equation}
\label{eq:xy-mother-bdg-block}
    \mathcal{H}_{\rm XY}(k)
    =
    \begin{pmatrix}
        \xi(k)&\Delta_\phi(k)\\
        \Delta_\phi(k)^*&-\xi(k)
    \end{pmatrix}.
\end{equation}
In the gauge $\phi=0$ and with \cref{eq:xy-mother-standard-gap}, this becomes
\begin{equation}
\label{eq:xy-mother-bdg-block-explicit}
    \mathcal{H}_{\rm XY}(k)
    =
    \begin{pmatrix}
        2(h-J\cos k)&2\ii J\gamma\sin k\\
        -2\ii J\gamma\sin k&-2(h-J\cos k)
    \end{pmatrix}.
\end{equation}
Equivalently,
\begin{equation}
\label{eq:xy-mother-d-vector}
    \mathcal{H}_{\rm XY}(k)
    =
    \bm d_{\rm XY}(k)\cdot\bm\tau,
    \qquad
    \bm d_{\rm XY}(k)
    =
    \bigl(0,-2J\gamma\sin k,2(h-J\cos k)\bigr),
\end{equation}
where $\tau_x,\tau_y,\tau_z$ act in Nambu space.  The sign of the $d_y$
component follows from the convention $d_y=-\operatorname{Im}\Delta(k)$.

The quasiparticle energies are
\begin{equation}
\label{eq:xy-mother-quasiparticle-spectrum}
    E_\pm(k)=\pm E(k),
    \qquad
    E(k)
    =
    2\sqrt{(h-J\cos k)^2+(J\gamma\sin k)^2}.
\end{equation}
This is the central exact result of the chapter.  All phase boundaries and low-energy limits follow from the zeros of this dispersion.

\begin{remark}[BdG eigenvalues versus many-body excitations]
The negative branch $-E(k)$ is the particle--hole partner of the positive branch.
It is not an independent physical particle.  After the Bogoliubov
transformation, the many-body Hamiltonian is written in terms of positive-energy
quasiparticles, with a ground-state energy shift.
\end{remark}

\section{Bogoliubov angle and diagonal form}
\label{sec:xy-bogoliubov-angle-diagonal-form}

For each independent non-self-conjugate pair $(k,-k)$, define
\begin{equation}
\label{eq:xy-mother-pair-xi-delta}
    \xi_k=2(h-J\cos k),
    \qquad
    \Delta_k=2\ii J\gamma\,\ee^{-2\ii\phi}\sin k,
    \qquad
    E_k=\sqrt{\xi_k^2+\abs{\Delta_k}^2}.
\end{equation}
The Bogoliubov polar angle is
\begin{equation}
\label{eq:xy-mother-bogoliubov-angle}
    \cos\theta_k=\frac{\xi_k}{E_k},
    \qquad
    \sin\theta_k=\frac{\abs{\Delta_k}}{E_k}.
\end{equation}
Writing
\begin{equation}
    \Delta_k=\abs{\Delta_k}\ee^{\ii\alpha_k},
\end{equation}
the coherence factors may be chosen as
\begin{equation}
\label{eq:xy-mother-coherence-factors}
    u_k=
    \sqrt{\frac{E_k+\xi_k}{2E_k}},
    \qquad
    v_k=
    -\ee^{\ii\alpha_k}
    \sqrt{\frac{E_k-\xi_k}{2E_k}}.
\end{equation}
The quasiparticle operators are
\begin{equation}
\label{eq:xy-mother-quasiparticle-operators}
    \gamma_k=u_kc_k-v_kc_{-k}^\dagger,
    \qquad
    \gamma_{-k}=u_kc_{-k}+v_kc_k^\dagger.
\end{equation}
They satisfy canonical anticommutation relations and diagonalize the pair block:
\begin{equation}
\label{eq:xy-mother-pair-diagonal}
    H_k
    =
    E_k
    \left(
        \gamma_k^\dagger\gamma_k
        +
        \gamma_{-k}^\dagger\gamma_{-k}
        -1
    \right).
\end{equation}
The normalized even-parity ground state of the pair block is
\begin{equation}
\label{eq:xy-mother-pair-ground-state}
    \ket{\Omega_k}
    =
    \left(
        u_k+v_k c_k^\dagger c_{-k}^\dagger
    \right)
    \ket{0}_{k,-k}.
\end{equation}
Thus, ignoring possible self-conjugate modes for the moment, the BCS ground state
in a fixed boundary sector is
\begin{equation}
\label{eq:xy-mother-bcs-ground-state}
    \ket{\Omega}
    =
    \prod_{k\in\mathcal{K}_+}
    \left(
        u_k+v_k c_k^\dagger c_{-k}^\dagger
    \right)
    \ket{0},
\end{equation}
where $\mathcal{K}_+$ contains one representative from each pair $\{k,-k\}$.

The corresponding ground-state energy in the thermodynamic limit is
\begin{equation}
\label{eq:xy-mother-ground-state-energy-density}
    e_0
    =
    \lim_{L\to\infty}\frac{E_0}{L}
    =
    -\frac{1}{2\pi}\int_0^\pi E(k)\,\dd k,
\end{equation}
with $E(k)$ given in \cref{eq:xy-mother-quasiparticle-spectrum}.  This formula
includes the constant $-hL$ in the original spin Hamiltonian.  For example, in
the large-field limit $h\gg J$, it gives $e_0\simeq -h$, as expected for the
fully $+z$ polarized spin state.

\section{Self-conjugate momenta and finite-size parity bookkeeping}
\label{sec:xy-self-conjugate-momenta}

The momenta $k=0$ and $k=\pi$ satisfy $k\equiv-k$ modulo $2\pi$.  They occur only
in the periodic fermion sector and only for compatible system sizes.  Since the
spinless pairing gap is odd in momentum,
\begin{equation}
    \Delta(0)=0,
    \qquad
    \Delta(\pi)=0,
\end{equation}
these modes are not paired.  They must be treated as ordinary fermionic modes.
Their normal energies are
\begin{equation}
\label{eq:xy-mother-self-conjugate-xi}
    \xi(0)=2(h-J),
    \qquad
    \xi(\pi)=2(h+J).
\end{equation}
Occupying such a mode changes fermion parity and therefore affects the exact
finite-size matching between the spin Hilbert space and the fermionic sector.

For bulk thermodynamics, the contribution of these isolated modes is negligible
as $L\to\infty$.  For exact diagonalization, finite-size spectra, and parity
selection, they are essential.  A safe finite-size procedure is:
\begin{enumerate}[label=(\roman*)]
    \item choose a spin boundary condition;
    \item split the fermionic problem into $P_F=+1$ and $P_F=-1$ sectors;
    \item impose APBC in the even sector and PBC in the odd sector;
    \item diagonalize each quadratic Hamiltonian;
    \item keep only states whose quasiparticle occupations have the required
    fermion parity.
\end{enumerate}

\section{Excitation gap and phase diagram}
\label{sec:xy-phase-diagram}

The bulk gap closes when
\begin{equation}
\label{eq:xy-mother-gap-closing-condition}
    h-J\cos k=0,
    \qquad
    J\gamma\sin k=0.
\end{equation}
Assume $J\neq0$.  For $\gamma\neq0$, the second equation forces
\begin{equation}
    k=0
    \quad\text{or}\quad
    k=\pi.
\end{equation}
The first equation then gives the two Ising critical lines
\begin{equation}
\label{eq:xy-mother-ising-critical-lines}
    h=J
    \quad(k=0),
    \qquad
    h=-J
    \quad(k=\pi).
\end{equation}
Near $h=J$, set $k\ll1$.  Then
\begin{equation}
\label{eq:xy-mother-near-h-plus-J}
    E(k)
    \simeq
    2\sqrt{
        \left(h-J+\frac{Jk^2}{2}\right)^2
        +(J\gamma k)^2
    }.
\end{equation}
At the critical point $h=J$ with $\gamma\neq0$, the leading dispersion is
linear,
\begin{equation}
\label{eq:xy-mother-ising-linear-dispersion}
    E(k)\simeq 2\abs{J\gamma}\,\abs{k},
\end{equation}
which is the lattice signature of an Ising/Majorana critical point with dynamical
exponent $z=1$.

Similarly, near $h=-J$ one expands around $k=\pi$.  Writing $q=k-\pi$, one finds
another Ising critical point with the same low-energy structure.

The situation changes qualitatively when $\gamma=0$.  The pairing term vanishes
and the Hamiltonian is number conserving.  The dispersion reduces to
\begin{equation}
\label{eq:xy-mother-xx-dispersion}
    \xi(k)=2(h-J\cos k).
\end{equation}
The system is gapless whenever there are Fermi points satisfying
\begin{equation}
\label{eq:xy-mother-fermi-points}
    \cos k_F=\frac{h}{J}.
\end{equation}
Thus, for $J>0$,
\begin{equation}
\label{eq:xy-mother-xx-critical-region}
    \gamma=0,
    \qquad
    \abs{h}<J
\end{equation}
defines the gapless XX region.  This is not an Ising critical line.  It is a
number-conserving free-fermion critical phase with two Fermi points and central
charge $c=1$ in the continuum limit.  By contrast, the Ising lines at
$\gamma\neq0$ have a single Majorana critical mode and central charge $c=1/2$.

The phase structure is therefore:
\begin{equation}
\label{eq:xy-mother-phase-summary}
\begin{array}{ccl}
\gamma\neq0,\ h=J &:& \text{Ising critical line at } k=0,\\[2pt]
\gamma\neq0,\ h=-J &:& \text{Ising critical line at } k=\pi,\\[2pt]
\gamma=0,\ \abs{h}<J &:& \text{gapless XX/free-fermion critical phase},\\[2pt]
\abs{h}>J &:& \text{gapped polarized phase},\\[2pt]
\abs{h}<J,\ \gamma\neq0 &:& \text{gapped paired/ordered phase}.
\end{array}
\end{equation}
At the multicritical points $(h,\gamma)=(\pm J,0)$, the linear pairing velocity
vanishes and the dispersion along the $\gamma=0$ direction becomes quadratic.
This is why the multicritical point is not described by the same scaling limit as
a generic Ising point.

\section{Relation to the transverse-field Ising model}
\label{sec:xy-relation-to-tfim}

Setting $\gamma=1$ in \cref{eq:xy-hamiltonian} gives
\begin{equation}
\label{eq:xy-mother-tfim-limit}
    H_{\rm TFIM}
    =
    -J\sum_j\sigma_j^x\sigma_{j+1}^x
    -h\sum_j\sigma_j^z.
\end{equation}
The quasiparticle spectrum becomes
\begin{equation}
\label{eq:xy-mother-tfim-spectrum}
    E_{\rm TFIM}(k)
    =
    2\sqrt{(h-J\cos k)^2+J^2\sin^2 k}
    =
    2\sqrt{h^2+J^2-2hJ\cos k}.
\end{equation}
For $J>0$ and $h\ge0$, the physically standard transition is at
\begin{equation}
\label{eq:xy-mother-tfim-critical-h-positive}
    h=J.
\end{equation}
The excitation gap is
\begin{equation}
\label{eq:xy-mother-tfim-gap}
    \Delta_{\rm gap}=2\abs{h-J}.
\end{equation}
The sign of $h-J$ becomes the sign of the continuum Majorana mass.  This is the
fermionic counterpart of the order-disorder transition of the Ising chain.

The second line $h=-J$ is present in the full XY parameter plane.  For the usual
TFIM convention one often restricts to $J>0$ and $h\ge0$, in which case this
second line is outside the standard positive-field discussion.  Keeping it in the
mother-model phase diagram is useful because it is required by the full lattice
BdG structure.

\section{Relation to the XX chain}
\label{sec:xy-relation-to-xx}

Setting $\gamma=0$ removes the pairing term in
\cref{eq:xy-mother-bulk-fermion-hamiltonian}.  The Hamiltonian becomes
\begin{equation}
\label{eq:xy-mother-xx-fermion-hamiltonian}
    H_{\rm XX}
    =
    -J\sum_j
    \left(
        c_j^\dagger c_{j+1}+c_{j+1}^\dagger c_j
    \right)
    -hL+2h\sum_j n_j.
\end{equation}
This is a tight-binding model with dispersion
\begin{equation}
\label{eq:xy-mother-xx-tight-binding-dispersion}
    \xi(k)=2(h-J\cos k).
\end{equation}
The ground state is a Fermi sea, not a BCS paired state.  For $J>0$, modes with
\begin{equation}
    \xi(k)<0
    \quad\Longleftrightarrow\quad
    \cos k>\frac{h}{J}
\end{equation}
are occupied.  When $\abs{h}<J$, the two Fermi points are
\begin{equation}
\label{eq:xy-mother-xx-fermi-momenta}
    k=\pm k_F,
    \qquad
    k_F=\arccos\left(\frac{h}{J}\right).
\end{equation}
Linearizing near the Fermi points gives two chiral fermion modes, or equivalently
a compact boson/Luttinger-liquid continuum theory.  This is why the XX critical
region has central charge $c=1$, not the Ising value $c=1/2$.

\section{Relation to the Kitaev chain}
\label{sec:xy-relation-to-kitaev-chain}

The bulk fermionic image of the XY chain is a spinless $p$-wave superconducting
chain.  To compare conventions cleanly, write the Kitaev chain as
\begin{equation}
\label{eq:xy-mother-kitaev-convention-annihilation}
    H_{\rm K}
    =
    -t\sum_j
    \left(
        c_j^\dagger c_{j+1}+c_{j+1}^\dagger c_j
    \right)
    -\mu\sum_j\left(n_j-\frac12\right)
    +\Delta\sum_j
    \left(
        c_jc_{j+1}+c_{j+1}^\dagger c_j^\dagger
    \right),
\end{equation}
with real $\Delta$ for the moment.  Since
\begin{equation}
    c_{j+1}^\dagger c_j^\dagger=-c_j^\dagger c_{j+1}^\dagger,
    \qquad
    c_jc_{j+1}=-c_{j+1}c_j,
\end{equation}
comparison with \cref{eq:xy-mother-bulk-fermion-expanded} gives the bulk
dictionary
\begin{equation}
\label{eq:xy-mother-kitaev-dictionary}
    t=J,
    \qquad
    \Delta=J\gamma,
    \qquad
    \mu=-2h.
\end{equation}
The constant term also matches:
\begin{equation}
    -\mu\sum_j\left(n_j-\frac12\right)
    =
    2h\sum_j n_j-hL
    \qquad
    \text{when } \mu=-2h.
\end{equation}

Some references instead define the pairing term as
\begin{equation}
    \Delta_{\rm cr}\sum_j
    \left(
        c_j^\dagger c_{j+1}^\dagger+c_{j+1}c_j
    \right).
\end{equation}
In that convention the same XY Hamiltonian corresponds to
\begin{equation}
\label{eq:xy-mother-kitaev-creation-dictionary}
    \Delta_{\rm cr}=-J\gamma.
\end{equation}
This is only a convention for how the Hermitian conjugate pair is written.  The
physical spectrum depends on $\abs{\Delta}$ and on the relative orientation of the
BdG vector, not on this bookkeeping sign alone.

Using \cref{eq:xy-mother-kitaev-dictionary}, the Kitaev-chain topological
condition
\begin{equation}
    \abs{\mu}<2\abs{t}
\end{equation}
becomes
\begin{equation}
\label{eq:xy-mother-kitaev-topological-condition-spin}
    \abs{h}<\abs{J}.
\end{equation}
This is precisely the region between the two Ising critical lines of the XY
mother model when $\gamma\neq0$.  In the spin language it is the ordered region
for the Ising-anisotropic chain; in the fermion language it is the topological
superconducting phase with Majorana edge modes under open boundary conditions.

\section{Topological interpretation of the BdG vector}
\label{sec:xy-bdg-vector-topological-interpretation}

For real $J$, $h$, and $\gamma$, and in a real pairing gauge, the BdG vector lies
in a two-dimensional plane:
\begin{equation}
\label{eq:xy-mother-bdg-plane-vector}
    \bm d(k)=
    \bigl(0,-2J\gamma\sin k,2(h-J\cos k)\bigr).
\end{equation}
Equivalently, one may encode the loop by the complex function
\begin{equation}
\label{eq:xy-mother-q-function}
    q(k)=d_z(k)+\ii d_y(k)
    =
    2(h-J\cos k)-2\ii J\gamma\sin k.
\end{equation}
As $k$ runs around the Brillouin zone, $q(k)$ traces an ellipse centered at
$2h$ on the real axis.  The origin is enclosed if and only if
\begin{equation}
    \abs{h}<\abs{J}
    \qquad
    \text{and}
    \qquad
    \gamma\neq0.
\end{equation}
The corresponding winding number may be written as
\begin{equation}
\label{eq:xy-mother-winding-number}
    \nu
    =
    \frac{1}{2\pi\ii}
    \int_{-\pi}^{\pi}\dd k\,\partial_k\log q(k).
\end{equation}
With the orientation convention in \cref{eq:xy-mother-q-function},
\begin{equation}
\label{eq:xy-mother-winding-result}
    \nu
    =
    \begin{cases}
        -\operatorname{sgn}(J\gamma), & \abs{h}<\abs{J},\\
        0, & \abs{h}>\abs{J}.
    \end{cases}
\end{equation}
Changing the sign convention for the pairing reverses the sign of $\nu$ but does
not change whether the phase is topological.  The robust class-D statement is
the corresponding $\ZZ_2$ distinction:
\begin{equation}
\label{eq:xy-mother-z2-topological-condition}
    \abs{h}<\abs{J}
    \quad\Longleftrightarrow\quad
    \text{nontrivial Kitaev phase}.
\end{equation}

\section{Continuum limits}
\label{sec:xy-continuum-limits}

The same exact lattice dispersion produces different continuum theories in
different regions of parameter space.

\subsection*{Ising/Majorana scaling near $h=J$}

Take $J>0$, $\gamma\neq0$, and expand near $k=0$ and $h=J$.  Keeping the leading
terms gives
\begin{equation}
\label{eq:xy-mother-majorana-expansion-h-plus}
    \mathcal{H}_{\rm XY}(k)
    \simeq
    2(h-J)\tau_z
    -2J\gamma k\,\tau_y,
\end{equation}
up to the sign convention of the pairing plane.  This is the momentum-space form
of a massive Majorana Hamiltonian.  The mass is proportional to
\begin{equation}
\label{eq:xy-mother-majorana-mass-h-plus}
    m_+=2(h-J),
\end{equation}
and the velocity is proportional to
\begin{equation}
    v=2J\gamma.
\end{equation}
The critical theory at $m_+=0$ has central charge $c=1/2$.

Near $h=-J$, one expands around $k=\pi$.  Setting $k=\pi+q$ gives another
Majorana theory with mass proportional to
\begin{equation}
\label{eq:xy-mother-majorana-mass-h-minus}
    m_-=2(h+J).
\end{equation}

\subsection*{XX scaling for $\gamma=0$}

For $\gamma=0$ and $\abs{h}<J$, expand around the two Fermi points
$\pm k_F$.  The Fermi velocity is
\begin{equation}
\label{eq:xy-mother-xx-fermi-velocity}
    v_F
    =
    \left.\abs{\frac{\dd \xi(k)}{\dd k}}\right|_{k=k_F}
    =
    2J\sin k_F
    =
    2J\sqrt{1-\left(\frac{h}{J}\right)^2}.
\end{equation}
The resulting low-energy theory is a massless Dirac fermion, or equivalently a
compact boson.  Adding a small nonzero $\gamma$ pairs the two Fermi points and
gaps the XX critical phase, except at the Ising critical endpoints.

\section{Summary}
\label{sec:xy-mother-summary}

The transverse-field XY chain is solved by the following sequence:
\begin{enumerate}[label=(\roman*)]
    \item Jordan--Wigner transformation maps the spin chain to a quadratic
    spinless fermion Hamiltonian with hopping, chemical potential, and $p$-wave
    pairing.
    \item Periodic spin boundary conditions split the fermionic problem into
    parity sectors: even parity gives APBC, odd parity gives PBC.
    \item Fourier transformation decomposes the Hamiltonian into independent
    $(k,-k)$ BdG blocks.
    \item Each block has
    \begin{equation}
        \xi(k)=2(h-J\cos k),
        \qquad
        \abs{\Delta(k)}=2\abs{J\gamma\sin k}.
    \end{equation}
    \item The exact quasiparticle dispersion is
    \begin{equation}
        E(k)=2\sqrt{(h-J\cos k)^2+(J\gamma\sin k)^2}.
    \end{equation}
    \item For $\gamma\neq0$, the bulk gap closes at $h=\pm J$, giving Ising
    critical lines.
    \item For $\gamma=0$, the chain is number conserving and is gapless for
    $\abs{h}<\abs{J}$, giving the XX/free-fermion critical phase.
    \item Under the Kitaev-chain dictionary $t=J$, $\mu=-2h$, and
    $\Delta=J\gamma$ up to pairing convention, the region $\abs{h}<\abs{J}$ is
    the topological superconducting phase.
\end{enumerate}

This single solution therefore supplies the exact solution of the TFIM, the XX
chain, and the Kitaev chain as special cases or equivalent descriptions.
